\section{Related Work}
\label{sec:related}

%Before we lay out our approach, it helps to see exactly what has already been tried on this problem, and where each attempt stops short.

\textit{Defects4J} \citep{just2014defects4j} is one of the most widely used benchmark datasets for software debugging and defect analysis research. It provides a curated collection of reproducible real-world bugs from 17 open-source Java projects, along with their original bug reports, reproducible failures, buggy and fixed source code, and buggy-to-fixed patches. This rich and reproducible set of artifacts makes \textit{Defects4J} particularly suitable for automated defect classification. In this research, we leverage these artifacts to investigate the classification of \textit{Defects4J} bugs into ODC defect categories using large language models (LLMs). The remainder of this section provides the necessary background on \textit{Defects4J}, ODC, and LLM-based defect classification.

\subsection{Existing \textit{Defects4J} Classification Attempts}

%Prior work on classifying \textit{Defects4J} bugs falls into three lines, each covering a different slice of the dataset with a different goal.

Van der Spuy et al. \citep{vanderspuy2025anatomy} manually classified 488 faults using control-flow and data-flow graph abstractions, identifying six control-flow and two data-flow categories. This structural perspective provides insight into how defects alter program execution; however, it does not necessarily capture the underlying reason for the defect. For example, a missing null check and an omitted feature may result in similar control-flow changes while representing fundamentally different defect causes.

Sobreira \textit{et al.} \citep{sobreira2018dissection} analyzed 395 bugs from \textit{Defects4J} v1.2, focusing on repair actions and patch patterns. Their findings showed that a substantial proportion of patches involve modifications to conditional blocks. Although patch-based analysis is valuable for understanding how developers repair defects, patch characteristics primarily describe the changes made to the program rather than the underlying defect mechanism that necessitated those changes.

Xuan \textit{et al.} \citep{xuan2017better} manually examined 50 randomly selected bugs to investigate test-suite quality and fault fixability. While their analysis provides detailed insights into individual defects, the relatively small sample limits the generalizability of the findings to the broader \textit{Defects4J} dataset.

Collectively, these studies provide valuable perspectives on defect structure, repair actions, and fixability. However, they examine relatively limited subsets of \textit{Defects4J}, rely primarily on code-structural or patch-based taxonomies, and do not apply a semantic defect-type taxonomy to the full set of bugs. Our work addresses this gap by investigating ODC-based semantic defect classification across the \textit{Defects4J} dataset using large language models.
\subsection{Orthogonal Defect Classification} \label{subsec:ODC}

ODC was introduced by Chillarege \textit{et al.}\ \citep{chillarege1992odc} as
a framework that classifies defects along orthogonal attributes. It was
originally designed to give an in-process signal during active development. This study applies the same attribute vocabulary
to defects already resolved in a static benchmark. These attributes are
recorded in two phases: \emph{Opener} attributes at discovery time and
\emph{Closer} attributes at fix time, summarised in
Table~\ref{tab:odc-attributes}.

\begin{table}[H]
\centering
\caption{ODC Opener and Closer attributes \citep{ibm2013odc}.}
\label{tab:odc-attributes}
\setlength{\tabcolsep}{5pt}
\begin{tabularx}{\columnwidth}{@{}llX@{}}
\toprule
\thead{Phase} & \thead{Attribute} & \thead{Description} \\
\midrule
\multirow{3}{*}{Opener} & Activity    & Specific activity being performed when the defect was found (e.g.\ code inspection or unit test) \\
                         & Trigger     & Condition that exposed the defect \\
                         & Impact      & User-visible consequence of the defect \\
\midrule
\multirow{5}{*}{Closer} & Defect Type & The actual correction made \\
                         & Target      & Design or Code \\
                         & Qualifier   & Missing, Incorrect, or Extraneous \\
                         & Age         & Whether the faulty code is base, new, rewritten, or refixed \\
                         & Source      & Whether the faulty code is developed in-house, reused from a library, outsourced, or ported \\
\bottomrule
\end{tabularx}
\end{table}

Of the eight attributes above, this study classifies Defect Type and Impact.
The Defect Type values for Design/Code targets are further organised into
seven types.

\subsection{\textcolor{red}{mapping}}
Going through the eight attributes from Table~\ref{tab:odc-attributes} in
order shows how much of each one Defects4J can actually support.

\begin{itemize}
\item \textbf{Activity.} Defects4J bugs are reproduced by running a unit
  test, so Unit Test is the natural fit for this attribute. Bug reports
  sometimes describe how the problem was noticed, which could point to a
  different activity such as inspection, but this shows up inconsistently
  across reports.
\item \textbf{Trigger.} The same reasoning applies here. The unit test that
  reproduces the bug is the closest available signal for the condition that
  exposed it, and a report's own description could suggest something else,
  again inconsistently.
\item \textbf{Impact.} This is what we derive directly from the bug report
  and the failure behavior, one of the two attributes this study actually
  classifies.
\item \textbf{Defect Type.} This is what the study tries to predict from the
  buggy code, then checks that prediction once the fixed code becomes
  available.
\item \textbf{Target.} By definition this is Code. Defects4J bugs are fixed
  in source, and there is no separate design document ODC's Design target
  would require.
\item \textbf{Qualifier.} This becomes available once the fix itself is
  visible, since it describes whether the fix supplied something missing,
  corrected something wrong, or removed something unneeded.
\item \textbf{Age.} Defects4J bugs come from real project histories with
  full commit logs, so age could plausibly be traced through the commit
  history of the affected code. This is not something pursued in this
  study.
\item \textbf{Source.} Most Defects4J projects are long running, community
  maintained Java libraries, so Source would intuitively be Developed
  In-House for nearly all of them. This is a reasonable expectation, though
  unconfirmed.
\end{itemize}

That leaves Defect Type and Impact as the two attributes this study
classifies. The others fall into three groups: fixed by how the dataset is
built, a reasonable guess without confirmation, or visible only once the
fix itself is known.

\subsection{LLM-based Defect Classification}

Give an LLM the right constraints and it classifies software artefacts
effectively. Colavito \textit{et al.}\ \citep{colavito2024llm} showed that LLMs
classify issue reports at an accuracy competitive with supervised ML, without
any labelled training data, and Koyuncu \citep{koyuncu2025exploring} extended
this to fine-grained bug-report categorisation using few-shot templates. Earlier
automated ODC approaches based on SVM and Naive Bayes reached only 77--82\%
accuracy and struggled with sparse text
\citep{huang2011autoodc,thung2012automatic}; LLMs sidestep these limits through
code-aware reasoning over code, tests, and error messages.

Chain-of-thought prompting improves classification by forcing structured
reasoning \citep{nong2024cot}. \textcolor{red}{Separately, Kang \textit{et al.}\
\citep{kang2023autosd} embedded Zeller's scientific debugging method
\citep{zeller2009why} directly into LLM prompts (AutoSD): the model generates a
hypothesis, runs a debugger command to test it, and uses the result to decide
whether to conclude, for explainable fault localization. We build on the same
iterative hypothesis--predict--experiment structure as AutoSD
\citep{kang2023autosd}, but repurpose it: our loop does not localize or repair
a fault, it \emph{classifies} the defect's ODC type, using probes into
held-back evidence as its experiments.} Gao \textit{et al.}\ \citep{gao2024rcegen} likewise showed that
LLMs can produce ODC-aligned root-cause explanations from reports and stack
traces given taxonomy guidance. Across these studies, taxonomy-free LLM
classification remains fragile: without a constraining label space, runs disagree
on the same bug \citep{koyuncu2025exploring}, and hallucinated labels are most
common when trace context is missing or ambiguous \citep{gao2024rcegen}.



\subsection{Research Gap}

The gap we address is the lack of a pipeline that is all four of these at once:
(i) \emph{context-rich}, grounded in multi-modal \textit{Defects4J} artefacts;
(ii) \emph{semantic}, classifying by defect type rather than syntactic
pattern; (iii) \emph{coverage-aware}, using an open-set ODC formulation that
measures rather than assumes taxonomy fit; and (iv) \emph{evaluable without
ground truth}, through paired pre-fix/post-fix self-consistency comparison. The
rest of this paper builds exactly that pipeline.

%% =============================================================================
